\subsection*{Morphisms between Maurer--Cartan forms}
% \label{morphisms}

Henceforth we drop the qualification ``generalized'': All Maurer--Cartan
forms and logarithmic derivatives will be understood in the
generalized sense.

With $G$, $M$ and $\Sigma $ fixed as in the Introduction (under
``Generalized Maurer--Cartan forms'') we collect all associated
Maurer--Cartan forms into the objects of a category. In~this category a
morphism $\omega_1 \rightarrow \omega_2$ between objects
$\omega_1 \colon A_1 \rightarrow {\mathfrak g} $ and
$\omega_2 \colon A_2 \rightarrow {\mathfrak g} $ consists of a Lie
algebroid morphism $\lambda \colon A_1 \rightarrow A_2$ covering the
identity on~$\Sigma $ and an element $l \in G$ such that the following
diagram commutes:
 \begin{equation*}
 %\label{morph}
 \begin{CD} A_1 @>{\omega_1}>> {\mathfrak g} \\ @V{\lambda}VV
@VV{\Adjoint_l}V \\ A_2 @>{\omega_2}>> {\mathfrak g}.
 \end{CD}
 \end{equation*} If $\lambda$ is injective, we will say that $\omega_1
\rightarrow \omega_2$ is {\it monic}. The preceding abstractions are
justified by the following observation (strengthened in special cases
in Theorem~\ref{uniquenessT} below):

\begin{Proposition}\label{morphismsP}
 Let $f_1\colon \Sigma \rightarrow M$ be a smooth map into a
 homogeneous $G$-space $M$ and define a second smooth map
 $f_2 \colon \Sigma \rightarrow M$ by $f_2 = \phi\circ f_1$, for
 some $\phi \in \automorphism(M)$. Then $\delta f_1$ and $\delta f_2$
 are isomorphic in the category of Maurer--Cartan forms.
\end{Proposition}

That is, smooth maps $f_1,f_2 \colon \Sigma \rightarrow M$ agreeing up
to a symmetry of $M$ have isomorphic logarithmic derivatives.

\begin{proof} Supposing $f_2 = \phi \circ f_1$, $\phi \in
\automorphism(M)$, define $l \in G$ as in Definition
\ref{symmetriesD}. Then the map $\Adjoint_l \times \phi $, defined by
 \begin{align*}
 (\xi,x) &\mapsto \big(\!\Adjoint_l \xi , \phi(x)\big),
 \\
{\mathfrak g} \times M &\rightarrow {\mathfrak g} \times M
 \end{align*}
 is a Lie algebroid automorphism of the action
algebroid ${\mathfrak g} \times M$ covering $\phi \colon M \rightarrow
M$. In~particular, the composite $A(f_1) \rightarrow {\mathfrak g}
\times M \xrightarrow{\Adjoint_l \times \phi} {\mathfrak g} \times M$
is a Lie algebroid morphism $J$ sitting in a~com\-mutative diagram
\begin{equation*}
 \begin{CD}
 A(f_1) @>{J}>> {\mathfrak g} \times M \\ @VVV @VVV \\
T \Sigma @>>{Tf_2}> TM.
 \end{CD}
\end{equation*}
The vertical arrows indicate anchor maps. Explicitly, we have
\[
J(X) = \big(\!\Adjoint_l \delta f_1 (X), f_2(\lrcorner X)\big),
\]
where $\lrcorner X \in \Sigma$ denotes the base point of $X$.

As $A(f_2)$ is the pullback of ${\mathfrak g} \times M$ under $f_2$,
we obtain, from the universal property of pullbacks, a unique Lie
algebroid morphism $\lambda \colon A(f_1) \rightarrow A(f_2)$ such
that $J $ is the composite
 \[
 A(f_1) \xrightarrow{\lambda} A(f_2) \xrightarrow{\delta f_2}
 {\mathfrak g} \times M.
 \]

This immediately implies commutativity of the diagram
\begin{equation*}
 \begin{CD} A(f_1) @>{\delta f_1}>> {\mathfrak g} \\
@V{\lambda}VV @VV{\Adjoint_l}V \\ A(f_2) @>{\delta f_2}>> {\mathfrak g}.
 \end{CD}
\end{equation*}
One argues that $\lambda$ is an isomorphism by
replacing $\phi $ with $\phi^{-1}$ and reversing the roles of $f_1$
and~$f_2$.
\end{proof}

\subsection*{Primitives}

A smooth map $f \colon \Sigma \rightarrow M$ will be called a {\it primitive} of the Maurer--Cartan form
$\omega \colon A \rightarrow {\mathfrak g} $ if there exists a
morphism $\omega \rightarrow \delta f$. Evidently, every principal
primitive is a primitive.

\subsection*{Maximal Maurer--Cartan forms}% \label{yef2}

Note that Axioms M1 and M2, together with Proposition \ref{frogP},
imply the following restrictions on the necessarily constant rank of
$A$, whenever $\omega \colon A \rightarrow {\mathfrak g} $ is a
Maurer--Cartan form:
\begin{equation*}
\dimension {\mathfrak g} -\dimension M\le \rank A
\le \dimension {\mathfrak g} -\dimension M + \dimension \Sigma.
\end{equation*}
We say $\omega $ is {\it maximal} if $A$ has maximal
rank, i.e., if
\[
\dim {\mathfrak g}-\rank A = \dimension M - \dimension \Sigma.
\]
In
this case it follows from M3, Proposition \ref{frogP}, and a dimension
count that
\begin{enumerate}\itemsep=0pt
\item[M3$'$.] For any point $x \in \Sigma $ there exists $m \in M$
such that $\omega(A^\circ_x) = {\mathfrak g}_m$.
\end{enumerate} Logarithmic derivatives and ordinary Maurer--Cartan
forms are always maximal.
%
\begin{Lemma}%\label{yef2L}
 Every morphism $\omega \rightarrow \delta f$ is monic. In~particular, if $\omega $ is maximal, then
 $\omega \rightarrow \delta f$ is an isomorphism.
\end{Lemma}

\begin{proof} A morphism $\omega \rightarrow \delta f$ consists of a
Lie algebroid morphism $\lambda \colon A \rightarrow A(f)$ covering
the identity on~$\Sigma$, and $l \in G$, such that
 \begin{equation}
 \delta f (\lambda(a))=\Adjoint_l \omega(a), \qquad a \in A.\label{cat}
 \end{equation}
 Suppose $\lambda(a)=0$, $a$ an element of $A $ with
base-point $x \in \Sigma $. Since $\lambda$ is a Lie algebroid
morphism covering the identity, we have $\# a =0$, i.e., $a \in
A^\circ_x$. Since $\omega(a)=0$, by \eqref{cat}, Axiom M2 and
Proposition \ref{frogP} imply $a=0$.
\end{proof}
%
The existence Theorem \ref{outlineT} has the following corollary (of
which we make no further use):
\begin{Corollary}%\label{yef2C}
 Every Maurer--Cartan form
 $\omega \colon A \rightarrow {\mathfrak g} $ with constant monodromy
 has an extension to a maximal Maurer--Cartan form
 $\omega \colon A' \rightarrow {\mathfrak g} $, for some Lie
 algebroid $A' \supset A$.
\end{Corollary}

\begin{proof} By the existence theorem, $\omega $ admits a principal
 primitive $f \colon\! \Sigma \rightarrow M$. That is, there exists a
 morphism $\lambda \colon\! A \rightarrow A(f)$, injective by the
 lemma, whose logarithmic derivative
 $\delta f \colon\! A(f) \rightarrow {\mathfrak g} $ fits into the
 commutative diagram~\eqref{carpet:court}. The logarithmic
 derivative of $f$ is then a maximal Maurer--Cartan form extending
 $\omega $.
\end{proof}

\subsection*{Uniqueness of primitives}% \label{uniqueness}

As usual, suppose $G$ acts transitively on $M$, and let
$G_{m_0}^\circ $ denote the connected component of the isotropy
$G_{m_0}$ at some $m_0 \in M$. Then since $G_{m_0}^\circ$ is {\em
 path}-connected, $N_G\big(G_{m_0}\big)\subset N_G\big(G_{m_0}^\circ\big)$.
\begin{Definition}%\label{uniquenessE}
 We say the isotropy groups of the $G$ action are {\it weakly
 connected} if for some (and hence any) $m_0 \in M$, we have
 $N_G\big(G_{m_0}\big)=N_G\big(G_{m_0}^\circ\big)$.
\end{Definition}

\begin{Example} If $M $ is one of the Riemannian space forms ${\mathbb
R}^n$, ${\mathbb S}^n$ or ${\mathbb H}^n $, and $G$ is the full group of
isometries, then although the isotropy groups of the action of $G$ on
$M$ are not connected, they {\em are} weakly connected.
\end{Example}
A proof of the following central result appears below.

\begin{Theorem}\label{uniquenessT}
 Suppose the action of $G$ on $M$ has weakly connected isotropy
 groups. Let $f_1, f_2 \colon \Sigma \rightarrow M$ be smooth
 maps. Then there exists an isomorphism $\delta f_1 \cong \delta f_2$
 in the category of Maurer--Cartan forms {\em if and only if} there
 exists $\phi \in \automorphism(M)$ such that
 $f_2 = \phi \circ f_1$.
\end{Theorem}

In contrast to the classical setting (Theorem~\ref{tslT})
there may exist more than one choice of $\phi \in \automorphism(M)$
for which $f_2 = \phi \circ f_1$, even if $G$ acts faithfully on
$M$. For example, consider two {\em constant} maps $f_1$, $f_2$.

Combining the theorem with the lemma above, we obtain:
\begin{Corollary}[uniqueness of primitives]\label{uniquenessC}
 If the action of $G$ on $M$ has weakly connected isotropy groups
 then primitives $f \colon \Sigma \rightarrow M$ of a maximal
 Maurer--Cartan form are unique, up to symmetries of $M$.
\end{Corollary}

A non-maximal Maurer--Cartan form may have distinct primitives not
related by a symmetry:
\begin{Example} Let $G$ be the group of isometries of the plane
$M={\mathbb R}^2 $ with Lie algebra ${\mathfrak g} $ identified with
the Killing fields. Let $x, y \colon {\mathbb R}^2 \rightarrow
{\mathbb R} $ denote the standard coordinate functions and let $
\omega \colon T {\mathbb R} \rightarrow {\mathfrak g} $ be the
generalized Maurer--Cartan form\footnote{Actually $\omega $ is an
ordinary Maurer--Cartan form in this case but we are understanding
primitives as maps into~${\mathbb R}^2$, not maps into $G$!} defined
by
\begin{equation*}
\omega \bigg( \frac{\partial }{\partial t}\bigg)=
-y\frac{\partial }{\partial x }+x \frac{\partial }{\partial y}
\qquad\text{(a constant element of $\mathfrak g$).}
 \end{equation*}
 Then for any $r\ge 0$ the map $f (t)=(r\cos t, r\sin
t)$ is a primitive of $\omega $.
\end{Example}

For the proof of the theorem we need one additional observation:

\begin{Proposition}\label{uniquenessP}
 Suppose $f \colon \Sigma \rightarrow M$ is a principal primitive of
 a Maurer--Cartan form $\omega \colon A \rightarrow {\mathfrak g} $.
 Let $\Omega \colon {\mathcal G} \rightarrow G$ be the global form of
 the monodromy of $\omega $, as defined in \eqref{global}. Then, for
 any $x \in \Sigma $, one has $x\xrightarrow{\omega}f(x)$, and for
 any $x_0\in \Sigma $,
 \begin{equation*} f(x)=\Omega(p)\cdot f(x_0),
 \end{equation*} where $p \in {\mathcal G} $ is any arrow from $x_0$
to $x$.
\end{Proposition}
\begin{proof} For some Lie algebroid morphism $\lambda \colon A
\rightarrow A(f)$, we have a commutative diagram
 \begin{equation*}
 \begin{CD} A @>{\omega }>> {\mathfrak g} \\ @V{\lambda}VV @AAA \\
A(f) @>>> {\mathfrak g} \times M.
 \end{CD}
 \end{equation*}
 Since $\lambda$ covers the identity, the claim
$x\xrightarrow{\omega}f(x)$ follows easily from commutativity and the
definition of the bottom map. Let ${\mathcal G}(f)$ denote the
pullback of the action groupoid $G \times M$ by $f$. Since ${\mathcal
G} $ is source-simply-connected, $\lambda$ is the derivative of a Lie
groupoid morphism $\Lambda \colon {\mathcal G} \rightarrow {\mathcal
G}(f)$ and the following diagram commutes (because the composites
being compared have a source-connected domain and identical
derivatives, by the commutativity of the preceding diagram):
 \begin{equation*}
 \begin{CD} {\mathcal G} @>{\Omega }>> G \\ @V{\Lambda}VV @AAA \\
{\mathcal G}(f) @>>> G\times M.
 \end{CD}
 \end{equation*}
 In particular, if we define $F$ to be the composite
Lie groupoid morphism ${\mathcal G}\xrightarrow{\Lambda}{\mathcal
G}(f)\rightarrow G \times M$, then $F$ covers $f \colon \Sigma
\rightarrow M$ and, by the commutativity,
\begin{equation}
F(p)=(\Omega(p),\alpha(p)),\label{rabbit}
\end{equation} where $\alpha $ denotes source projection. But as $F
\colon {\mathcal G} \rightarrow G \times M$ must respect the target
projections, denoted $\beta$, we also have
$f(\beta(p))=\beta(F(p))$. Now \eqref{rabbit} gives
\begin{equation*} f(\beta(p))=\Omega(p) \cdot \alpha(p),
\end{equation*} which proves the proposition.
\end{proof}
%
\begin{proof}[Proof of theorem (for $\boldsymbol{\Sigma}$ simply-connected)] That
$\delta f_1$ and $\delta f_2$ must be isomorphic when $f_2 = \phi
\circ f_1 $, $\phi \in \automorphism(M)$, is Proposition
\ref{morphismsP}. Suppose $\delta f_1 \cong\delta f_2$ and assume
initially that $\Sigma $ is simply-connected (needed in the proof of
the lemma below). By definition, there exists $l \in G$ and a Lie
algebroid isomorphism $\lambda \colon A(f_2) \rightarrow A(f_1)$ such
that the following diagram commutes:
\begin{gather}
 \begin{CD} A(f_1) @>{\delta f_1}>> {\mathfrak g} \\ @A{\lambda}AA
@VV{\Adjoint_l}V \\ A(f_2) @>{\delta f_2}>> {\mathfrak g}.\label{cod}
 \end{CD}
\end{gather}
Arbitrarily fixing a point $x_0 \in \Sigma$,
\eqref{uht} gives
\begin{equation}
\delta f_1\big(A(f_1)_{x_0}\big) = {\mathfrak g}_{f_1(x_0)},\qquad
\delta f_2\big(A(f_2)_{x_0}\big) = {\mathfrak g}_{f_2(x_0)}.\label{hws}
\end{equation}

For $i=1$ or $2$, let $\Omega_i \colon {\mathcal G}^i \rightarrow G$
denote the global form of the monodromy of $\delta f_i$, as defined at
\eqref{global}. The Lie algebroid of ${\mathcal G}^i$ is $A(f_i)$ and,
by Lie II for Lie groupoids, there is a unique Lie groupoid
isomorphism $\Lambda \colon {\mathcal G}^2\rightarrow {\mathcal G}^1$
whose derivative is $\lambda$. Taking $\omega = \delta f_i$ in the
preceding proposition, we obtain
\begin{equation} f_1(x) = \Omega_1(p_1) \cdot f_1(x_0),\qquad f_2(x) =
\Omega_2(p_2) \cdot f_2(x_0), \label{pto}
\end{equation} whenever $p_i \in {\mathcal G}^i$ is an arrow from
$x_0$ to $x$. By the commutativity of \eqref{cod}, the Lie groupoid
morphisms $\Omega_2 \colon {\mathcal G}^2 \rightarrow G$ and $p_2
\mapsto l \Omega_1({\Lambda}(p_2))l^{-1}$ have the same
derivative, namely $\delta f_2$, so they must coincide, because
${\mathcal G}^2$ is source-connected:
\begin{equation}
\Omega_2(p_2) = l\Omega_1({\Lambda}(p_2))l^{-1},\qquad
p_2 \in {\mathcal G}^2.\label{ksd}
\end{equation}
Since $\lambda$, and hence $\Lambda$, covers the
identity on~$\Sigma $, $p_1 \in {\mathcal G}^1$ is an arrow from $x_0$
to $x$ if and only if $p_2:=\Lambda(p_1)\in {\mathcal G}^2$ is an
arrow from $x_0$ to $x$. This fact and \eqref{ksd} allow us to rewrite
the second equation in \eqref{pto} as $f_2(x)=l \Omega_1(p_1)l^{-1}
\cdot f_2(x_0)$. Or, choosing $r \in G$ such that
\begin{equation} f_2(x_0)=lr^{-1}\cdot f_1(x_0),\label{five}
\end{equation} we have
\begin{equation} f_1(x) = \Omega_1(p_1) \cdot f_1(x_0),\qquad f_2(x) =
l\Omega_1(p_1)r^{-1} \cdot f_1(x_0), \label{pto2}
\end{equation} whenever $p_1 \in {\mathcal G}^1$ is an arrow from
$x_0$ to $x$.
\begin{Lemma} $r \in G$ lies in the normaliser of $G_{f_1(x_0)}$.
\end{Lemma}

Assuming the lemma holds, there exists, by the
characterization of symmetries in Proposition~\ref{symmetriesP}, an
element $\phi \in \automorphism(M)$ well-defined by $\phi(g \cdot
f_1(x_0))=lgr^{-1}\cdot f_1(x_0)$. Then \eqref{pto2} gives us
$f_2(x)=\phi(f_1(x))$, as required.
\end{proof}
%
\begin{proof}[Proof of lemma] Since we assume the isotropy groups of
the action of $G$ on $M$ are weakly connected, it suffices to show $r
\in N_G\big(G_{f_1(x_0)}^\circ\big)$. We claim
 \begin{gather}
 \Omega_1\big({\mathcal G}^1_{x_0}\big)=G_{f_1(x_0)}^\circ ,\label{hws2}
 \\
\Omega_2\big({\mathcal G}^2_{x_0}\big)=G_{f_2(x_0)}^\circ.\label{hws3}
 \end{gather}
 Since ${\mathcal G}^i$ is transitive and
source-connected ($i \in \{1,2\}$) the restriction of the target
projection of ${\mathcal G}^i$ to the source-fibre over $x_0$ is a
principal ${\mathcal G}^i_{x_0}$-bundle over $\Sigma $. Since we
assume $\Sigma $ is simply-connected, ${\mathcal G}^i_{x_0}$ is
connected, by the long exact homotopy sequence for this principal
bundle. It~follows that \eqref{hws2} and \eqref{hws3} are consequences
of their infinitesimal analogues, which already appear in \eqref{hws}
above.

 Because $\Lambda \colon {\mathcal G}^2 \rightarrow {\mathcal G}^1$
is a Lie groupoid isomorphism covering the identity, we have
 \begin{equation}
 \Lambda \big({\mathcal G}^2_{x_0}\big) ={\mathcal G}^1_{x_0}.\label{kk}
 \end{equation}
 We now compute
 \begin{align*}
 r G_{f_1(x_0)}^\circ r^{-1} &= l^{-1}G_{lr \cdot
f_1(x_0)}^\circ l = l^{-1}G_{f_2(x_0)}^\circ l=l^{-1} \Omega_2
\big({\mathcal G}^2_{x_0}\big)l = \Omega_1
\big(\Lambda\big({\mathcal G}^2_{x_0}\big)\big)
\\
&=\Omega_1 \big({\mathcal
G}^1_{x_0}\big)=G_{f_1(x_0)}^\circ.
 \end{align*}
 The second and subsequent equalities in this
computation follow from equations \eqref{five}, \eqref{hws3},
\eqref{ksd}, \eqref{kk} and \eqref{hws2} respectively.
\end{proof}

\begin{proof}[Proof of theorem {\normalfont(general case)}]\sloppy If $\delta f_1 \cong
\delta f_2$ but $\Sigma $ is not simply-connected, then \mbox{$\delta(f_1
\circ \pi)\cong \delta(f_2 \circ \pi)$}, where $\pi \colon \tilde
\Sigma \rightarrow \Sigma $ denotes the universal covering map, as it
is not difficult to see. By the result just proven in the
simply-connected case, there exists $\phi \in \automorphism(M)$ such
that $f_1 \circ \pi = \phi \circ f_2 \circ \pi $. But as $\pi $ is
surjective, this immediately implies $f_1 = \phi \circ f_2 $.
\end{proof}

\subsection*{Summary of results}
% \label{summary}

Suppose $f$ is a primitive of a Maurer--Cartan form $\omega $, so that
$\delta f (\lambda(X))=\Adjoint_l \omega(X)$, for some Lie algebroid
morphism $\lambda $ and element $l \in G$. Then it is not hard to show
that $f'(x)=l^{-1} \cdot f(x)$ defines a {\em principal} primitive of
$\omega $. That is, the existence of primitives already implies the
existence of principal primitives. We may therefore summarise the
results cited in the Introduction and our uniqueness result, Corollary
\ref{uniquenessC}, as follows:

\begin{Theorem}[main theorem]\label{summaryT}
 Let $M$ be a homogeneous $G$-space and
 $\omega \colon A \rightarrow {\mathfrak g} $ an associated
 generalized Maurer--Cartan form, where $A$ is a Lie algebroid over
 some manifold $\Sigma $. Then~$A$ is integrable. Furthermore,
 $\omega $ admits a primitive $f \colon \Sigma \rightarrow M$ if and
 only if it has constant monodromy
 $\Omega_{x_0}^{m_0} \colon \pi_1(\Sigma, x_0) \rightarrow M$, for
 some choice of $x_0 \in \Sigma $ and $m_0 \in M$ with
 $x_0\xrightarrow{\omega}m_0$. Assu\-ming~$\omega $ is maximal, and the isotropy groups of the action of $G$ on $M$ are weakly
 connected, the primitive $f$ is unique up to symmetry.
\end{Theorem}

We reiterate that ``symmetry'' is to be understood in the sense Definition \ref{symmetriesD}.

